\documentclass[10pt,a4paper]{amsart}
\usepackage[margin=19mm]{geometry}
\usepackage{amsmath,amssymb,mathtools}
\usepackage[scr=rsfs,cal=euler]{mathalfa}
\usepackage{xfrac}
\usepackage[unicode,hidelinks,bookmarks=false]{hyperref}
\newcommand{\ie}{i.e.\ }
\newcommand{\cf}{cf.\ }
\newcommand{\resp}{respectively\ }
\newcommand{\Subsection}{\subsection}
\providecommand{\nobreakdash}{\nobreak\hbox{-}}
\setlength{\emergencystretch}{2em}
\multlinegap0pt
\usepackage{mathrsfs}
\newcommand{\Ric}{\operatorname{Ric}}
\newcommand{\Hess}{\operatorname{Hess}}
\newcommand{\grad}{\operatorname{grad}}
\newcommand{\divg}{\operatorname{div}}
\newcommand{\Vol}{\operatorname{Vol}}
\newcommand{\tr}{\operatorname{tr}}
\newtheorem{theorem}{Theorem}[section]
\newtheorem{lemma}[theorem]{Lemma}
\newtheorem{proposition}[theorem]{Proposition}
\newtheorem{corollary}[theorem]{Corollary}
\newtheorem{example}[theorem]{Example}
\newtheorem{remark}[theorem]{Remark}
\newtheorem{definition}[theorem]{Definition}
\begin{document}
\title[Rigidity of complete non-compact generalized m-quasi-Einstei]{Rigidity of complete non-compact generalized m-quasi-Einstein manifolds}
\author{M. Ahmad Mirshafeazadeh}
\date{}
\maketitle
\noindent\emph{Source and licence.} Rigidity of complete non-compact generalized m-quasi-Einstein manifolds. Version arXiv:2606.06129v2 (CC BY 4.0). This material is distributed under the Creative Commons Attribution 4.0 International licence, http://creativecommons.org/licenses/by/4.0/, as recorded in the arXiv OAI licence metadata for this exact version and shown on the abstract page https://arxiv.org/abs/2606.06129v2. Full rights notice: licenses/arxiv-2606.06129-CC-BY-4.0.txt.\par\smallskip
\noindent\textit{Editorial scope.} Counted unit: the original \texttt{theorem} environment \texttt{thm4} at source lines 520--532 together with its complete proof at lines 534--571; the delivered document keeps source lines 251--572 so that every referenced label and every citation resolves inside it. Native editable source is the paper's own concatenated TeX; the AMS wrapper is editorial formatting only. No publisher class, upstream script or unrelated artwork is redistributed.
\begin{equation}
R_{ij} + f_{ij} - \mu f_i f_j = \lambda g_{ij}, \label{eq:GQE2}
\end{equation}
where $f_{ij} = \nabla_j\nabla_i f$ and $f_i = \nabla_i f$.
When $\mu$ is the constant $1/m$ for a positive real number $m$, we speak of a gradient generalized $m$-quasi-Einstein manifold.

From \eqref{eq:GQE2} we immediately obtain the trace
\begin{equation}
\Delta f = n\lambda - R + \mu|\nabla f|^2. \label{eq:trace}
\end{equation}
Contracting \eqref{eq:GQE2} with $f^j$ yields a useful gradient formula:
\begin{equation}
\frac12\nabla_i|\nabla f|^2 = \lambda f_i - R_{ij}f^j + \mu|\nabla f|^2 f_i. \label{eq:grad}
\end{equation}

\subsection{The scalar quantities $I$ and $J$}

We recall the two auxiliary functions introduced in \cite{Mir14} that govern the Laplacian of the scalar curvature.
For a gradient GQE manifold we define
\[
\begin{aligned}
I :=&\; (n-1)\Delta\lambda - 2\mu(n-1)\langle\grad f,\grad\lambda\rangle + \frac{1+2\mu}{2}\langle\grad R,\grad f\rangle \\
&\quad - (1-\mu)|\hat{\Ric}|^2 + \frac{(n\lambda-R)\bigl(R(1-\mu+\mu n) + \mu\lambda n(1-n)\bigr)}{n},
\end{aligned}
\]
where $\hat{\Ric} = \Ric - \frac{R}{n}g$ is the trace-free Ricci tensor, and a similar expression for $J$ (which we do not need explicitly because it will vanish under our assumptions).
The next identity was proved in \cite{Mir14} (see also \cite{Mir15}).
\begin{lemma}\label{lem:I}
On any gradient GQE manifold,
\begin{equation}
\frac12\Delta R = I + J. \label{eq:DeltaR}
\end{equation}
\end{lemma}

In the present paper we will restrict to the case where both the scalar curvature $R$ and the coefficient $\mu$ are constant.
Under these assumptions all derivatives of $R$ and $\mu$ vanish; consequently $J\equiv 0$ and $\Delta R=0$, so Lemma~\ref{lem:I} forces $I=0$.
For $\mu=1/m$ this condition reads
\begin{equation}
\begin{aligned}
0 = (n-1)\Delta\lambda &- \frac{2}{m}(n-1)\langle\grad f,\grad\lambda\rangle - \frac{m-1}{m}|\hat{\Ric}|^2 \\
&+ \frac{(n\lambda-R)\bigl(R(1+\frac{n-1}{m}) + \frac{\lambda n(1-n)}{m}\bigr)}{n}. \label{eq:Izero}
\end{aligned}
\end{equation}
Equation \eqref{eq:Izero} will be our starting point for the derivation of the Laplacian of $v$.

\subsection{Analytical toolbox}
\label{sec:toolbox}

We collect the classical results that will be used in the proofs.
All manifolds are assumed to be complete and without boundary.

\begin{lemma}[Yau \cite{Yau}]\label{lem:Yau}
Let $u$ be a non-negative smooth subharmonic function on a complete Riemannian manifold $M$.
If $u\in L^p(M)$ for some $p>1$, then $u$ is constant.
\end{lemma}

\begin{lemma}[Schoen--Yau \cite{SchoenYau}]\label{lem:SY}
Let $u\ge 0$ be a smooth subharmonic function on a complete manifold $M$.
For any ball $B(q,2r)\subset M$ and any smooth cut-off function $\phi$ supported in $B(q,2r)$ with $\phi\equiv 1$ on $B(q,r)$ and $|\nabla\phi|\le c/r$, there exists a universal constant $A$ (depending only on $c$) such that
\[
\int_{B(q,r)}|\nabla u|^2 \le \frac{A}{r^2}\int_{B(q,2r)} u^2.
\]
\end{lemma}

\begin{lemma}[Caminha--Souza--Camargo \cite{Caminha}]\label{lem:CSC}
Let $X$ be a smooth vector field on a complete, non-compact, oriented Riemannian manifold $M$.
If $\divg X$ does not change sign on $M$ and $|X|\in L^1(M)$, then $\divg X\equiv 0$.
\end{lemma}
\begin{remark}
If $M$ is not orientable, we may pass to the orientation double cover $\widetilde M$.
Every structure (metric, soliton functions, etc.) lifts isometrically, the lifted vector field $\widetilde X$ satisfies the same hypotheses, and the conclusion $\divg \widetilde X=0$ on $\widetilde M$ implies $\divg X=0$ on $M$ by local isometry.
Hence the lemma applies without loss of generality.
\end{remark}

\begin{lemma}[Karp \cite{Karp}]\label{lem:Karp}
Let $u\ge 0$ be a smooth subharmonic function on a complete, non-compact Riemannian manifold with non-negative sectional curvature.
If $\int_M |\nabla u|^{\frac{n}{n-1}} < \infty$, then $u$ is constant.
\end{lemma}

\begin{lemma}[Tashiro \cite{Tashiro}]\label{lem:Tashiro}
Let $(M,g)$ be a complete Riemannian manifold.
If there exists a smooth function $w$ on $M$ such that $\Hess w = c\,g$ for some constant $c\neq 0$, then $M$ is isometric to Euclidean space, and $w$ is a quadratic polynomial of the form
\[
w(x) = \frac{c}{2}|x|^2 + \langle a, x\rangle + b.
\]
\end{lemma}

Finally, we need the existence of good cut-off functions under a mild Ricci lower bound.
The following statement is a special case of results of Bianchi and Setti \cite{Bianchi}.

\begin{lemma}[{\cite[Corollary 2.3]{Bianchi}}]\label{lem:Bianchi}
Assume that the Ricci curvature of a complete manifold $M$ satisfies
\[
\Ric \ge -(n-1)\frac{k^2}{1+r^2}\, g
\]
for some constant $k$, where $r(x)=d(x,q)$ is the distance from a fixed point $q\in M$.
Then, for every large $r>0$, there exists a smooth cut-off function $\xi_r$ with compact support in $B(q,2r)$ such that
\[
0\le \xi_r\le 1,\quad \xi_r\equiv 1 \text{ on } B(q,r),\quad
|\nabla\xi_r|\le \frac{K}{r},\quad |\Delta\xi_r|\le \frac{K}{r^2},
\]
with a constant $K>0$ independent of $r$.
\end{lemma}

\section{The key identity – subharmonicity of $v$}
\label{sec:key}

In this section we assume that $M$ is a gradient generalized $m$-quasi-Einstein manifold with constant scalar curvature $R$ and $\mu=1/m$, and we set
\[
v = e^{-f/m}\lambda.
\]

\begin{lemma}\label{lem:key}
Under the above assumptions,
\begin{equation}
e^{f/m}\Delta v = \frac{1}{n-1}\left(\frac{m-1}{m}|\hat{\Ric}|^2
- \frac{R(n\lambda-R)}{n}\Bigl(1+\frac{n-1}{m}\Bigr)\right). \label{eq:key}
\end{equation}
\end{lemma}

\begin{proof}
From \eqref{eq:Izero} we solve for $\Delta\lambda$:
\begin{align}
\Delta\lambda &= \frac{2}{m}\langle\grad f,\grad\lambda\rangle
+ \frac{m-1}{m(n-1)}|\hat{\Ric}|^2 \notag\\
&\quad - \frac{(n\lambda-R)\bigl(R(1+\frac{n-1}{m}) + \frac{\lambda n(1-n)}{m}\bigr)}{n(n-1)}. \label{eq:Deltalambda}
\end{align}
On the other hand, a direct computation using $\divg(\phi X)=\phi\divg X+\langle\grad\phi,X\rangle$ gives
\begin{align*}
e^{f/m}\Delta v &= e^{f/m}\divg\bigl( e^{-f/m}(\nabla\lambda - \tfrac{\lambda}{m}\nabla f) \bigr)\\
&= \Delta\lambda - \frac{2}{m}\langle\grad f,\nabla\lambda\rangle
- \frac{1}{m}\lambda\Delta f + \frac{1}{m^2}\lambda|\nabla f|^2.
\end{align*}
Now use the trace identity \eqref{eq:trace}, $\Delta f = n\lambda - R + \frac{1}{m}|\nabla f|^2$, to replace $\Delta f$:
\[
e^{f/m}\Delta v = \Delta\lambda - \frac{2}{m}\langle\grad f,\nabla\lambda\rangle
- \frac{n}{m}\lambda^2 + \frac{1}{m}\lambda R.
\]
Insert the expression for $\Delta\lambda$ from \eqref{eq:Deltalambda}; the gradient terms cancel and we obtain
\[
e^{f/m}\Delta v = \frac{m-1}{m(n-1)}|\hat{\Ric}|^2
- \frac{(n\lambda-R)\bigl(R(1+\frac{n-1}{m}) + \frac{\lambda n(1-n)}{m}\bigr)}{n(n-1)}
- \frac{\lambda}{m}(n\lambda-R).
\]
The last two summands combine as follows.
Let us compute
\begin{align*}
&-\frac{(n\lambda-R)}{n(n-1)}\Bigl(R(1+\tfrac{n-1}{m}) + \frac{\lambda n(1-n)}{m}\Bigr) - \frac{\lambda}{m}(n\lambda-R)\\
&= -(n\lambda-R)\Bigl[ \frac{R(1+\frac{n-1}{m})}{n(n-1)} + \frac{\lambda n(1-n)}{m}\cdot\frac{1}{n(n-1)} + \frac{\lambda}{m} \Bigr].
\end{align*}
Since $\frac{\lambda n(1-n)}{m}\cdot\frac{1}{n(n-1)} = \frac{\lambda(1-n)}{m(n-1)} = -\frac{\lambda}{m}$, the bracket simplifies to
\[
\frac{R(1+\frac{n-1}{m})}{n(n-1)} - \frac{\lambda}{m} + \frac{\lambda}{m}
= \frac{R(1+\frac{n-1}{m})}{n(n-1)}.
\]
Thus the whole expression becomes
\[
-(n\lambda-R)\frac{R(1+\frac{n-1}{m})}{n(n-1)}
= -\frac{R(n\lambda-R)}{n(n-1)}\Bigl(1+\frac{n-1}{m}\Bigr).
\]
Adding the first summand, we obtain exactly \eqref{eq:key}.
\end{proof}

When $R\le 0$, $\lambda>0$ and $m>1$, every term on the right-hand side of \eqref{eq:key} is clearly non-negative; hence $\Delta v\ge 0$.
Since $v=e^{-f/m}\lambda>0$, $v$ is a positive subharmonic function.

\section{Proof of Theorem 1 ($v\in L^p$)}
\label{sec:thm1}

\begin{theorem}\label{thm1}
Let $(M^n,g)$ be a complete non-compact gradient generalized $m$-quasi-Einstein manifold with constant scalar curvature $R\le 0$, $\lambda>0$ and $m>1$.
If the weighted function $v=e^{-f/m}\lambda$ belongs to $L^p(M)$ for some $p>1$, then \textbf{no such manifold exists}. Equivalently, the hypotheses are mutually inconsistent.
\end{theorem}

\begin{proof}
By Lemma~\ref{lem:key}, $v>0$ is subharmonic.
Yau's lemma (Lemma~\ref{lem:Yau}) implies that $v$ is constant; write $v\equiv c>0$.
Then $\Delta v=0$, and the non-negative right-hand side of \eqref{eq:key} must vanish term by term:
\begin{gather*}
\frac{m-1}{m(n-1)}|\hat{\Ric}|^2 = 0,\qquad
-\frac{R(n\lambda-R)}{n(n-1)}\Bigl(1+\frac{n-1}{m}\Bigr) = 0.
\end{gather*}
Since $m>1$, the first equality forces $|\hat{\Ric}|^2=0$, i.e.\ $\Ric = \frac{R}{n}g$.
The second, together with $n\lambda-R>0$ (because $\lambda>0$ and $R\le 0$), gives $R=0$.
Hence $\Ric = 0$.

The GQE equation reduces to $\Hess f - \frac1m df\otimes df = \lambda g$, and $v=c$ yields $\lambda = c e^{f/m}$.
Set $w = m e^{-f/m}$.
A direct computation gives
\[
\Hess w = \nabla( -\nabla f\, e^{-f/m}) = e^{-f/m}\bigl( -\Hess f + \frac1m df\otimes df \bigr)
= -e^{-f/m}\lambda g = -c\,g.
\]
Thus $w$ is a smooth positive function (since $w = m e^{-f/m} > 0$) with $\Hess w = -c\,g$, where $c>0$.
By Tashiro's theorem (Lemma~\ref{lem:Tashiro}), $M$ is isometric to Euclidean space, and $w$ must be of the form
\[
w(x) = -\frac{c}{2}|x|^2 + \langle a, x\rangle + b
\]
for some $a\in\mathbb{R}^n$, $b\in\mathbb{R}$.
But this expression is strictly negative for sufficiently large $|x|$, contradicting $w>0$ everywhere.
Therefore, the assumptions lead to a contradiction; hence no such complete non-compact manifold exists.
\end{proof}

\section{Proof of Theorem 2 (linear volume growth and bounded $v$)}
\label{sec:thm2}

\begin{theorem}\label{thm2}
Let $(M^n,g)$ be a complete non-compact gradient generalized $m$-quasi-Einstein manifold with constant scalar curvature $R\le 0$, $\lambda>0$ and $m>1$.
Assume that $v$ is bounded above and that the volume of geodesic balls grows at most linearly, i.e.\ $\Vol(B(q,r)) \le C r$ for some $C>0$ and all $r>0$.
Then \textbf{no such manifold exists}.
\end{theorem}

\begin{proof}
Let $v\le C_0$ on $M$.
By Lemma~\ref{lem:SY} applied to the subharmonic function $v$, for any large $r$ we have
\[
\int_{B(q,r)}|\nabla v|^2 \le \frac{16}{r^2}\int_{B(q,2r)} v^2
\le \frac{16 C_0^2}{r^2}\Vol(B(q,2r)).
\]
Using the linear volume growth, $\Vol(B(q,2r))\le 2C r$, we obtain
\[
\int_{B(q,r)}|\nabla v|^2 \le \frac{32 C_0^2 C}{r},
\]
which tends to $0$ as $r\to\infty$.
Thus $\int_M |\nabla v|^2 = 0$, so $\nabla v=0$ everywhere; consequently $v$ is constant.
The remainder of the proof is identical to that of Theorem~\ref{thm1}, leading to the same contradiction with $w>0$.
\end{proof}

\section{Proof of Theorem 3 ($|\nabla v|\in L^1$)}
\label{sec:thm3}

\begin{theorem}\label{thm3}
Let $(M^n,g)$ be a complete non-compact gradient generalized $m$-quasi-Einstein manifold with constant scalar curvature $R\le 0$, $\lambda>0$ and $m>1$.
If $|\nabla v|\in L^1(M)$, then \textbf{no such manifold exists}.
\end{theorem}

\begin{proof}
Apply Lemma~\ref{lem:CSC} to the vector field $X = \nabla v$ (after passing to the orientation cover if necessary, see Remark after Lemma~\ref{lem:CSC}).
From Lemma~\ref{lem:key} we have $\divg X = \Delta v \ge 0$, and by assumption $|X| = |\nabla v|\in L^1(M)$.
Thus $\Delta v \equiv 0$.
Vanishing of the non-negative right-hand side of \eqref{eq:key} gives $\hat{\Ric}=0$ and $R=0$, so $\Ric=0$.

The GQE equation becomes $\Hess f - \frac1m df\otimes df = \lambda g$. 
We now prove that $v$ must be constant without invoking Kanai's lemma.
Set $w = m e^{-f/m}$. Then, as before,
\[
\Hess w = -v\,g.
\]
Taking the trace gives $\Delta w = -n v$.
Now apply the divergence operator to both sides of $\Hess w = -v\,g$:
\[
\divg(\Hess w) = -\nabla v.
\]
On the other hand, in Riemannian geometry we have the standard identity
\[
\divg(\Hess w) = \nabla(\Delta w) + \Ric(\nabla w, \cdot).
\]
Since $\Ric=0$, this reduces to
\[
\divg(\Hess w) = \nabla(\Delta w) = \nabla(-n v) = -n\nabla v.
\]
Comparing with $-\nabla v$, we obtain $(n-1)\nabla v = 0$, and since $n\ge 3$, $\nabla v = 0$. Hence $v$ is constant.
Now $v\equiv c>0$ and the same argument as in Theorem~\ref{thm1} (using Tashiro's theorem and the positivity of $w$) yields a contradiction.
Therefore no such manifold exists.
\end{proof}

\section{Proof of Theorem 4 (weighted $L^1$ under a Ricci lower bound)}
\label{sec:thm4}


\begin{theorem}\label{thm4}
Let $(M^n,g)$ be a complete non-compact gradient generalized $m$-quasi-Einstein manifold with constant scalar curvature $R\le 0$, $\lambda>0$ and $m>1$.
Suppose that there exists a constant $k$ such that
\[
\Ric \ge -(n-1)\frac{k^2}{1+r(x)^2}\,g,
\]
where $r(x)=d(x,q)$ for a fixed point $q\in M$, and that the function $v$ satisfies the weighted integrability condition
\[
\int_{M\setminus B(q,r)} \frac{v(x)}{d(x,q)^2}\,dV_g(x) < \infty
\]
for all $r>0$ sufficiently large.
Then \textbf{no such manifold exists}.
\end{theorem}

\begin{proof}
Let $\Phi = e^{f/m}\Delta v$; by Lemma~\ref{lem:key}, $\Phi\ge 0$ and $\Delta v = e^{-f/m}\Phi$.
Thanks to the Ricci lower bound, Lemma~\ref{lem:Bianchi} supplies cut-off functions $\xi_r$ with the stated properties.
Multiply the identity $\Delta v = e^{-f/m}\Phi$ by $\xi_r$ and integrate over $M$:
\[
\int_M \xi_r e^{-f/m}\Phi \,dV_g = \int_M \xi_r \Delta v \,dV_g.
\]
Using Green's second identity and the fact that $\xi_r$ is compactly supported, this equals $\int_M v \Delta \xi_r \,dV_g$.
Thus
\[
0 \le \int_M \xi_r e^{-f/m}\Phi = \int_M v \Delta \xi_r
= \int_{B(q,2r)\setminus B(q,r)} v \Delta \xi_r,
\]
since $\xi_r\equiv 1$ on $B(q,r)$ and $\xi_r=0$ outside $B(q,2r)$.
Now $|\Delta \xi_r|\le K/r^2$ and $v>0$, so
\[
\int_{B(q,2r)\setminus B(q,r)} v \Delta \xi_r \le \frac{K}{r^2}\int_{B(q,2r)\setminus B(q,r)} v.
\]
In the annular region $B(q,2r)\setminus B(q,r)$ we have $d(x,q)\le 2r$, hence
$\frac{1}{r^2} = \frac{4}{(2r)^2} \le \frac{4}{d(x,q)^2}$.
Consequently,
\[
\frac{K}{r^2}\int_{B(q,2r)\setminus B(q,r)} v
\le 4K \int_{B(q,2r)\setminus B(q,r)} \frac{v}{d(x,q)^2}.
\]
By hypothesis, the right-hand side tends to $0$ as $r\to\infty$ (the integral over the complement of a ball goes to zero).
Therefore
\[
\lim_{r\to\infty}\int_M \xi_r e^{-f/m}\Phi = 0.
\]
Because $\xi_r\uparrow 1$ pointwise and the integrand is non-negative, the monotone convergence theorem yields
\[
\int_M e^{-f/m}\Phi = 0.
\]
Consequently $\Phi\equiv 0$ almost everywhere, and by smoothness $\Phi\equiv 0$ everywhere.
Hence $|\hat{\Ric}|^2=0$ and $R=0$, so $\Ric=0$.
The rest of the proof is identical to that of Theorem~\ref{thm3}, leading to the same contradiction.
\end{proof}


\begin{thebibliography}{99}
\bibitem[Bianchi]{Bianchi} D. Bianchi, A. Setti, \textit{Laplacian cut-offs, porous and fast diffusion on manifolds and other applications}, Calc. Var. Partial Differential Equations \textbf{57} (2018), no. 1, Paper No. 4.
\bibitem[Caminha]{Caminha} A. Caminha, P. Souza, F. Camargo, \textit{Complete foliations of space forms by hypersurfaces}, Bull. Braz. Math. Soc. \textbf{41} (2010), 339--353.
\bibitem[Karp]{Karp} L. Karp, \textit{On Stokes' theorem for noncompact manifolds}, Proc. Amer. Math. Soc. \textbf{82} (1981), 487--490.
\bibitem[Mir14]{Mir14} M. A. Mirshafeazadeh, B. Bidabad, \textit{On generalized quasi-Einstein manifolds}, Adv. Pure Appl. Math. \textbf{10} (2019), no. 3, 193--202. \url{https://doi.org/10.1515/apam-2017-0112}.
\bibitem[Mir15]{Mir15} M. A. Mirshafeazadeh, B. Bidabad, \textit{On the rigidity of generalized quasi-Einstein manifolds}, Bull. Malays. Math. Sci. Soc. (2) \textbf{43} (2020), 2029--2042. \url{https://doi.org/10.1007/s40840-019-00788-8}.
\bibitem[SchoenYau]{SchoenYau} R. Schoen, S.-T. Yau, \textit{Lectures on Differential Geometry}, International Press, Cambridge, MA, 1994.
\bibitem[Tashiro]{Tashiro} Y. Tashiro, \textit{Complete Riemannian manifolds and some vector fields}, Trans. Amer. Math. Soc. \textbf{117} (1965), 251--275.
\bibitem[Yau]{Yau} S.-T. Yau, \textit{Some function-theoretic properties of complete Riemannian manifolds and their applications to geometry}, Indiana Univ. Math. J. \textbf{25} (1976), 659--670.
\end{thebibliography}
\end{document}
